\BourbakiExercisesSection

\begin{enumerate}
\def\labelenumi{\arabic{enumi}.}
\tightlist
\item
  If A and B are K-algebras and \(\alpha\) and \(\beta\) two K-homomorphisms of A into B, an \((\alpha, \beta)\) -derivation of A into B is a K-linear mapping d of A into B satisfying the relation
\end{enumerate}

\[
d (x y) = (d x) \alpha (y) + \beta (x) (d y) \quad \text { for } \quad x, y \quad \text { in } \quad A,
\]

in other words, a derivation in the sense of Definition 1 (no. 2) where \(\mathbf{A} = \mathbf{A}' = \mathbf{A}''\) , \(\mathbf{B} = \mathbf{B}' = \mathbf{B}''\) , \(d = d' = d''\) , \(\mu: \mathbf{A} \times \mathbf{A} \to \mathbf{A}\) is multiplication, \(\lambda_1: \mathbf{B} \times \mathbf{A} \to \mathbf{B}\) is the K-bilinear mapping \((z, x) \mapsto z\alpha(x)\) and \(\lambda_2: \mathbf{A} \times \mathbf{B} \to \mathbf{B}\) the K-bilinear mapping \((x, z) \mapsto \beta(x)z\) .

Suppose henceforth that A and B are commutative fields and that \(\alpha \neq \beta\) . Show that there then exists an element \(b \in \mathbf{B}\) such that \(d\) is the \((\alpha, \beta)\) -derivation \(x \mapsto b\alpha(x) - \beta(x)b\) . (It can be assumed that \(d \neq 0\) . Show first that the kernel N of \(d\) is the set of \(x \in \mathbf{A}\) such that \(\alpha(x) = \beta(x)\) and then show that for \(x \notin \mathbf{N}\) the element \(b = d(x)/( \alpha(x) - \beta(x))\) is independent of the choice of \(x\) .)

\begin{enumerate}
\def\labelenumi{\arabic{enumi}.}
\setcounter{enumi}{1}
\tightlist
\item
  Let K be a commutative field of characteristic \(\neq2\) , A and B two K-algebras and \(\alpha\) a K-linear mapping of A into B; an \(\alpha\) -derivation of A into B is a K-linear mapping d of A into B satisfying the relation
\end{enumerate}

\[
d (x y) = (d x) \alpha (y) + \alpha (x) (d y) \quad \text { for } \quad x, y \quad \text { in } \quad A,
\]

in other words, a derivation on the sense of Definition 1 with \(\mathbf{A} = \mathbf{A}' = \mathbf{A}''\) , \(\mathbf{B} = \mathbf{B}' = \mathbf{B}''\) , \(d = d' = d''\) , \(\mu\) multiplication in \(\mathbf{A}\) , \(\lambda_1\) the mapping \((z,x) \mapsto z\alpha(x)\) of \(\mathbf{B} \times \mathbf{A}\) into \(\mathbf{B}\) and \(\lambda_2\) the mapping \((x,z) \mapsto \alpha(x)z\) of \(\mathbf{A} \times \mathbf{B}\) into \(\mathbf{B}\) .

Show that if B is a commutative field, an extension of K, and \(d \neq 0\) , there exists \(b \neq 0\) in B such that

\[
\alpha (x y) = \alpha (x) \alpha (y) + b (d x) (d y) \quad \text { for } \quad x, y \quad \text { in } \quad A.
\]

If \(b = c^2\) with \(c \in \mathbf{B}\) , \(c \neq 0\) , there exist two homomorphisms \(f, g\) of \(\mathbf{A}\) into \(\mathbf{B}\) such that \(\alpha(x) = (f(x) + g(x)) / 2\) , \(dx = (f(x) - g(x)) / 2c\) . Otherwise, there exists a quadratic algebra (§ 2, no. 3) \(\mathbf{B}'\) over \(\mathbf{B}\) , a \(\mu \in \mathbf{B}'\) such that \(\mu^2 = c\) , and a homomorphism \(f\) of \(\mathbf{A}\) into \(\mathbf{B}'\) such that \(\alpha(x) = (f(x) + \overline{f(x)}) / 2\) , \(dx = (f(x) - \overline{f(x)}) / 2\) .

¶ 3. Let K be a field, commutative or otherwise, E the left vector space \(\mathrm{K}_{s}^{(\mathbf{N})}\) and let \((e_{n})_{n\in\mathbf{N}}\) be the canonical basis of this space. On the other hand let \(\sigma\) be an endomorphism of K and d a \((\sigma,1_{\mathbf{K}})\) -derivation (Exercise 1) of K into itself, in other words an additive mapping of K into itself such that

\[
d (\xi \eta) = (d \xi) \sigma (\eta) + \xi (d \eta).
\]

The elements of E are linear combinations of products \(\alpha_{n}e_{n}\) and a multiplication is defined on E (Z-bilinear mapping of E × E into E) by the conditions

\[
\begin{array}{l} (\alpha e _ {n}) (\beta e _ {m}) = (\alpha e _ {n - 1}) (\sigma (\beta) e _ {m + 1} + (d \beta) e _ {m}) \text { for } n \geqslant 1 \\ (\alpha e _ {0}) (\beta e _ {m}) = (\alpha \beta) e _ {m}. \end{array}
\]

Thus an (in general non-commutative) ring structure is defined on E with \(e_{0}\) as unit element (identified with the element l of K). If we write \(X = e_{1}\) , then \(e_{n} = X^{n}\) for \(n \geqslant 1\) , so that the elements of E can be written as

\[
\alpha_ {0} + \alpha_ {1} X + \dots + \alpha_ {n} X ^ {n},
\]

with the multiplication rule

\[
\mathrm{X} \alpha = \sigma (\alpha) \mathrm{X} + (d \alpha) \quad \text { for } \quad \alpha \in \mathrm{K}.
\]

This ring is denoted by K{[}X; σ, d{]} and is called the ring of non-commutative polynomials in X with coefficients in K, relative to σ and d.~For a non-zero element \(z = \sum_{j} \alpha_{j} X^{j}\) of E, the degree deg(z) is the greatest integer j such that \(\alpha_{j} \neq 0\) .

\begin{enumerate}
\def\labelenumi{(\alph{enumi})}
\tightlist
\item
  Show that if \(u, v\) are two elements \(\neq 0\) of \(E\) , then:
\end{enumerate}

\[
\begin{array}{r l} u v \neq 0 & \text { and } \quad \deg (u v) = \deg (u) + \deg (v); \\ \deg (u + v) & \leqslant \sup (\deg (u), \deg (v)) \quad \text { if } u + v \neq 0. \end{array}
\]

Moreover, if \(\deg(u) \geqslant \deg(v)\) , there exist two elements w, r of E such that \(u = wv + r\) and, either r = 0, or \(\deg(r) < \deg(v)\) (``Euclidean division'' in E).

\begin{enumerate}
\def\labelenumi{(\alph{enumi})}
\setcounter{enumi}{1}
\tightlist
\item
  Conversely, let R be a ring such that there exists a mapping \(u \mapsto \deg(u)\) of the set \(R'\) of elements \(\neq 0\) of R into N, satisfying the conditions of (a). Show that the set K consisting of 0 and the non-zero elements \(u \in R\) such that \(\deg(u) = 0\) is a (not necessarily commutative) field. If \(K \neq R\) and \(x\) is an element of R for which the integer \(\deg(x) > 0\) is the smallest possible, show
\end{enumerate}

that every element of R can be written uniquely in the form \(\sum_{i} \alpha_j x^j\) with \(\alpha_j \in K\) . (To see that every element of R is of this form, argue by reductio ad absurdum by considering an element \(u \in R\) which is not of this form and for which \(\deg(u)\) is the smallest possible.) Prove that there exist an endomorphism \(\sigma\) of K and a \((\sigma, 1_K)\) -derivation \(d\) of K such that \(x\alpha = \sigma(\alpha)x + d\alpha\) for all \(\alpha \in K\) . Deduce that R is isomorphic to a field or a ring K{[}X; \(\sigma, d]\) .

\begin{enumerate}
\def\labelenumi{\arabic{enumi}.}
\setcounter{enumi}{3}
\tightlist
\item
  Let M be a graded K-module of type N; show that every graded endomorphism of degree r can be extended uniquely to a derivation of degree r of the universal anticommutative algebra \(\mathbf{G}(\mathbf{M})\) (§ 7, Exercise 13).
\end{enumerate}

* 5. Let K be a field of characteristic \(p > 0\) , B the field K(X) of rational functions in one indeterminate and A the subfield K(X \(^{p}\) ) of B. Show that the canonical homomorphism \(\Omega_{\mathbf{K}}(\mathbf{A}) \otimes_{\mathbf{A}} \mathbf{B} \to \Omega_{\mathbf{K}}(\mathbf{B})\) is not injective.*

